\documentclass[11pt]{article}
\usepackage{mathblatt}
\def\mitzone{1}
\begin{document}
\blattkopf{Flächeninhalt und Volumen im Raum}{Gesamt}
\verzeichniszeile{\verz{zone}{Kennst du schon (Nr.~1–8)} \verztrenn \verz{e1}{1 Dreiecksflächen (Nr.~9–11)} \verztrenn \verz{e2}{2 Vierecksflächen (Nr.~12–14)} \verztrenn \verz{e3}{3 Einheit 3 (Nr.~15–16)} \verztrenn \verz{e4}{4 Einheit 4 (Nr.~17–18)} \verztrenn \verz{abhaken}{Das kann ich}}
% Zone „Das kennst du schon“ – aus bank/flaecheninhalt-und-volumen-im-raum/zone.jsonl (zusammenbau v0.1)
\einheitenkopf[zone][Kennst du schon]{Das kennst du schon}
\setcounter{aufgabe}{0}

%% TODO Titel: Ich-kann-Satz fehlt in der Bank (Platzhalter: Kettenname) – Nr. 1
%% TODO Anweisung: ein Satz über der ersten Teilaufgabe fehlt in der Bank – Nr. 1
\begin{aufgabe}{Flächenformeln der ebenen Figuren (Dreieck, Trapez, Parallelogramm, Raute/Drachen, Kreis) kennen und anwenden}
\begin{teile}
\teil Dreieck: Grundseite 8\,cm, Höhe dazu 5\,cm. Flächeninhalt? \leerfeld[cm²]
\teil Trapez: parallele Seiten 6{,}5\,cm und 3{,}5\,cm, Höhe 4\,cm. Flächeninhalt? \leerfeld[cm²]
\end{teile}
\end{aufgabe}

%% TODO Titel: Ich-kann-Satz fehlt in der Bank (Platzhalter: Kettenname) – Nr. 2
%% TODO Anweisung: ein Satz über der ersten Teilaufgabe fehlt in der Bank – Nr. 2
\begin{aufgabe}{Volumen- und Oberflächenformeln der Elementarkörper (Prisma, Pyramide, Zylinder, Kegel) kennen}
\begin{teile}
\teil Quader: 5\,cm lang, 4\,cm breit, 3\,cm hoch. Volumen? \leerfeld[cm³]
\teil Zylinder: Radius 2\,cm, Höhe 10\,cm. Volumen? Runde auf eine Stelle. \leerfeld[cm³]
\end{teile}
\end{aufgabe}

%% TODO Titel: Ich-kann-Satz fehlt in der Bank (Platzhalter: Kettenname) – Nr. 3
%% TODO Anweisung: ein Satz über der ersten Teilaufgabe fehlt in der Bank – Nr. 3
\begin{aufgabe}{Verbindungsvektoren bilden, Beträge berechnen, Mittelpunkte bestimmen}
\begin{teile}
\teil Verbindungsvektor von $P(1 | 2 | 3)$ nach $Q(4 | 6 | 3)$? \leerfeld
\teil Mittelpunkt der Strecke von $P(-2 | 5 | 1)$ nach $Q(4 | -1 | 6)$? M( \leerfeld | \leerfeld | \leerfeld )
\end{teile}
\end{aufgabe}

%% TODO Titel: Ich-kann-Satz fehlt in der Bank (Platzhalter: Kettenname) – Nr. 4
%% TODO Anweisung: ein Satz über der ersten Teilaufgabe fehlt in der Bank – Nr. 4
\begin{aufgabe}{Skalarprodukt und „null heißt senkrecht“}
\begin{teile}
\teil Skalarprodukt $(2 | 1 | 0) \cdot (1 | -2 | 5)$? \leerfeld
\teil Stehen $(3 | -2 | 1)$ und $(2 | 4 | 2)$ senkrecht aufeinander? \leerfeld
\end{teile}
\end{aufgabe}

%% TODO Titel: Ich-kann-Satz fehlt in der Bank (Platzhalter: Kettenname) – Nr. 5
%% TODO Anweisung: ein Satz über der ersten Teilaufgabe fehlt in der Bank – Nr. 5
\begin{aufgabe}{Satz des Pythagoras für Seitenhöhen und fehlende Katheten, Strahlensätze und Ähnlichkeit für Streckfaktoren}
\begin{teile}
\teil Rechtwinkliges Dreieck, Katheten 6\,cm und 8\,cm. Hypotenuse? \leerfeld[cm]
\teil Rechtwinkliges Dreieck, Hypotenuse 2{,}5\,m, eine Kathete 1{,}5\,m. Andere Kathete? \leerfeld[m]
\end{teile}
\end{aufgabe}

%% TODO Titel: Ich-kann-Satz fehlt in der Bank (Platzhalter: Kettenname) – Nr. 6
%% TODO Anweisung: ein Satz über der ersten Teilaufgabe fehlt in der Bank – Nr. 6
\begin{aufgabe}{Gleichungen aus Formeln lösen (nach einer Größe umstellen, Wurzelgleichungen, quadratische Terme, Prozentrechnung für Anteile)}
\begin{gleichungsraster}[2]
\gl{\text{Löse $\tfrac{1}{2} \cdot 6 \cdot h = 18$ nach $h$.}} &
\gl{\text{Löse $\tfrac{1}{3} \cdot 27 \cdot h = 45$ nach $h$.}} \\
\end{gleichungsraster}
\end{aufgabe}

%% TODO Titel: Ich-kann-Satz fehlt in der Bank (Platzhalter: Kettenname) – Nr. 7
%% TODO Anweisung: ein Satz über der ersten Teilaufgabe fehlt in der Bank – Nr. 7
\begin{aufgabe}{Flächenformeln der ebenen Figuren (Dreieck, Trapez, Parallelogramm, Raute/Drachen, Kreis) kennen und anwenden – Fehler finden}
\begin{teile}
\teil Ein Dreieck hat die Grundseite 10\,cm, eine schräge Seite 5\,cm und die Höhe 4\,cm zur Grundseite. Lisa rechnet: \rechnung{A &= \tfrac{1}{2} \cdot 10 \cdot 5 \\ &= 25\,\mathrm{cm}^2} Finde den Fehler und rechne richtig.
\end{teile}
\end{aufgabe}

%% TODO Titel: Ich-kann-Satz fehlt in der Bank (Platzhalter: Kettenname) – Nr. 8
%% TODO Anweisung: ein Satz über der ersten Teilaufgabe fehlt in der Bank – Nr. 8
\begin{aufgabe}{Flächenformeln der ebenen Figuren (Dreieck, Trapez, Parallelogramm, Raute/Drachen, Kreis) kennen und anwenden – selbst rechnen}
\begin{teile}
\teil Dreieck: Grundseite 12\,cm, schräge Seite 10\,cm, Höhe zur Grundseite 8\,cm. Flächeninhalt? \leerfeld[cm²]
\end{teile}
\end{aufgabe}

\clearpage
% Einheit 1 von 4 – Dreiecksflächen (bank/flaecheninhalt-und-volumen-im-raum/e1.jsonl, zusammenbau v0.1)
%% TODO Zeitmarke Einheit 1: Marken-Zeile nicht lesbar
%% TODO Zweigzeile Teil 1: was hier gelernt wird (Satz in Schülersprache aus dem Einheitstitel) – Einheit 1
\einheitenkopf[e1]{Einheit 1 von 4 · Dreiecksflächen}
\zweigzeile{Abitur GK}
\setcounter{aufgabe}{8}

%% TODO Titel: Ich-kann-Satz fehlt in der Bank (Platzhalter: Kettenname) – Nr. 9
%% TODO Anweisung: ein Satz über der ersten Teilaufgabe fehlt in der Bank – Nr. 9
\begin{aufgabe}{Welche Figur, welche Formel?}
\begin{teile}
\teil Ein Dach hat vier gleiche dreieckige Seitenflächen. Welche Formel für eine Seitenfläche? \\ \kreuz{Grundseite mal Höhe} \\ \kreuz{halbe Grundseite mal Höhe} \\ \kreuz{halbes Diagonalenprodukt}
\end{teile}
\end{aufgabe}

%% TODO Titel: Ich-kann-Satz fehlt in der Bank (Platzhalter: Kettenname) – Nr. 10
%% TODO Anweisung: ein Satz über der ersten Teilaufgabe fehlt in der Bank – Nr. 10
\begin{aufgabe}{Dreiecksflächen}
%% TODO \rechenplatz{n}: Schrittzahl des Grundfalls fehlt in der Bank (nur bei fester Schreibform, 2.3 b)
\begin{teile}
\teil Im Dreieck $PQR$ ist $M$ der Mittelpunkt von $PQ$, und $MR$ steht senkrecht auf $PQ$. Was ist $MR$? \\ \kreuz{Höhe zu PQ} \\ \kreuz{schräge Seite}
\teil $P(1 | 1 | 0)$, $Q(4 | 5 | 0)$, $R(-4 | 11 | 0)$, rechtwinklig bei $Q$. Flächeninhalt des Dreiecks $PQR$? A = \leerfeld FE
\teil Das Dreieck mit $P(1 | 0 | 2)$, $Q(5 | 4 | 0)$, $R(5 | 1 | 3)$ ist gleichschenklig mit der Basis $PQ$. Flächeninhalt? A = \leerfeld FE
\teil Die Ebene $3x - 2y + z = 12$ schneidet die $x$-Achse und die $y$-Achse. Diese beiden Schnittpunkte und der Ursprung bilden ein Dreieck. Flächeninhalt? A = \leerfeld FE
\teil Licht fällt in Richtung $(0 | 1 | -1)$ auf das Dreieck mit $P(1 | 0 | 2)$, $Q(4 | 0 | 2)$, $R(1 | 1 | 4)$. Bestimme die Schattenpunkte in der $x$-$y$-Ebene, zeichne das Schattendreieck ein und berechne seinen Flächeninhalt.

\begin{ksys3}[x1min=-1,x1max=6,x2min=-1,x2max=6,x3min=0,x3max=5] \rpunkt{1}{0}{2}{P} \rpunkt{4}{0}{2}{Q} \rpunkt{1}{1}{4}{R} \end{ksys3}
A = \leerfeld FE
\end{teile}
\begin{gleichungsraster}[2]
\gl{\text{Der Ursprung $O$, $P(6 | 8 | 0)$ und $R(3 | 4 | z)$ mit $z > 0$ bilden ein gleichschenkliges Dreieck mit der Basis $OP$ und dem Flächeninhalt 40. Wert von $z$? (Abitur 2022 GK)}} &
\end{gleichungsraster}
\begin{teile}
\teil $P(1 | 0 | 0)$, $Q(3 | 1 | 0)$, $R(2 | 4 | 1)$; für $T$ gilt $\overrightarrow{OT} = \overrightarrow{OR} - 3 \cdot \overrightarrow{PQ}$. Bestimme $T$ und das Verhältnis der Flächeninhalte von Dreieck $PQR$ und Viereck $PQRT$. \hfill (Abitur 2024 GK)
\teil Für die Punkte $R$, $S$, $T$ gilt $\overrightarrow{OT} = \overrightarrow{OR} + 2 \cdot \overrightarrow{RS}$. Begründe ohne Rechnung anhand einer Skizze, dass die Dreiecke $ORS$ und $OST$ denselben Flächeninhalt haben. \hfill (Abitur 2023 LK)
\end{teile}
\end{aufgabe}

%% TODO Titel: Ich-kann-Satz fehlt in der Bank (Platzhalter: Kettenname) – Nr. 11
%% TODO Anweisung: ein Satz über der ersten Teilaufgabe fehlt in der Bank – Nr. 11
\begin{aufgabe}{Dreiecksflächen – Fehler finden, Begründen, Darstellungswechsel, Anwendung}
\begin{teile}
\teil Das Dreieck mit $P(0 | 0 | 1)$, $Q(6 | 0 | 1)$, $R(3 | 4 | 1)$ ist gleichschenklig mit der Basis $PQ$. Tom rechnet: \rechnung{|PR| &= 5 \\ A &= \tfrac{1}{2} \cdot 6 \cdot 5 = 15} Finde den Fehler und rechne richtig.
\teil Begründe: Im gleichschenkligen Dreieck geht die Höhe zur Basis durch den Mittelpunkt der Basis.
\teil Stelle für das Dreieck $PQR$ im Koordinatensystem einen Flächenterm „Rechteck minus Randdreiecke“ auf und berechne den Flächeninhalt.

\begin{ksys}[xmin=-1,xmax=6,ymin=-1,ymax=5] \punkt{0}{1}{P} \punkt{5}{0}{Q} \punkt{3}{4}{R} \end{ksys}
\teil Ein dreieckiges Sonnensegel ist an zwei Pfosten in $P(0 | 0 | 2)$ und $Q(6 | 0 | 2)$ und an der Hauswand in $R(3 | 4 | 5)$ befestigt (1 LE = 1 m); es ist gleichschenklig mit der Basis $PQ$. Der Stoff kostet 12\,€ pro m². Was kostet das Segel? \leerfeld[€]
\end{teile}
\end{aufgabe}

\clearpage
% Einheit 2 von 4 – Vierecksflächen (bank/flaecheninhalt-und-volumen-im-raum/e2.jsonl, zusammenbau v0.1)
%% TODO Zeitmarke Einheit 2: Marken-Zeile nicht lesbar
%% TODO Zweigzeile Teil 1: was hier gelernt wird (Satz in Schülersprache aus dem Einheitstitel) – Einheit 2
\einheitenkopf[e2]{Einheit 2 von 4 · Vierecksflächen}
\zweigzeile{Abitur GK}
\setcounter{aufgabe}{11}

%% TODO Titel: Ich-kann-Satz fehlt in der Bank (Platzhalter: Kettenname) – Nr. 12
%% TODO Anweisung: ein Satz über der ersten Teilaufgabe fehlt in der Bank – Nr. 12
\begin{aufgabe}{Vierecksflächen}
%% TODO \rechenplatz{n}: Schrittzahl des Grundfalls fehlt in der Bank (nur bei fester Schreibform, 2.3 b)
\begin{teile}
\teil Ein Viereck hat vier gleich lange Seiten, seine Diagonalen sind verschieden lang. Welche Formel für den Flächeninhalt? \\ \kreuz{Seite mal Seite} \\ \kreuz{halbe Summe der parallelen Seiten mal Höhe} \\ \kreuz{halbes Diagonalenprodukt}
\teil $P(1 | 0 | 0)$, $Q(1 | 4 | 3)$, $R(1 | 10 | -5)$. Zeige, dass bei $Q$ ein rechter Winkel liegt, und berechne den Flächeninhalt des Rechtecks $PQRS$. A = \leerfeld FE
\teil Das Viereck mit $P(3 | 4 | 4)$, $Q(-3 | 7 | 4)$, $R(-1 | 2 | 0)$, $S(5 | -1 | 0)$ ist eine Raute. Flächeninhalt? A = \leerfeld FE
\teil Das Trapez $PQRS$ mit $P(0 | 1 | 1)$, $Q(8 | 1 | 1)$, $R(6 | 4 | 5)$, $S(2 | 4 | 5)$ hat die parallelen Seiten $PQ$ und $SR$ und gleich lange Schenkel. Flächeninhalt? \hfill (Abitur 2019 GK) \\ A = \leerfeld FE
\teil Ein Quadrat liegt in der $x$-$y$-Ebene, eine Ecke ist der Ursprung $O$. Sein Diagonalenschnittpunkt $M$ liegt auf der Geraden $\vec x = (2 | 4 | 4) + t \cdot (2 | 2 | -2)$. Bestimme $M$ und den Flächeninhalt des Quadrats. \hfill (Abitur 2024 GK)
\end{teile}
\begin{gleichungsraster}[2]
\gl{\text{Das Rechteck $PQRS$ mit $P(1 | 2 | 0)$, $Q(4 | 6 | 0)$ und $S(1 | 2 | k)$, $k > 0$, hat den Flächeninhalt 30. Wert von $k$?}} &
\end{gleichungsraster}
\begin{teile}
\teil Das Quadrat $PQRS$ mit $P(0 | 1 | 0)$, $Q(8 | 1 | 0)$, $R(8 | 9 | 0)$, $S(0 | 9 | 0)$ wird verändert: $P$ und $S$ werden auf der $y$-Achse gleich weit aufeinander zu verschoben, sodass ein gleichschenkliges Trapez $P'QRS'$ mit drei Vierteln der Quadratfläche entsteht. Bestimme $P'$ und $S'$. \hfill (Abitur 2021 GK)
\end{teile}
\end{aufgabe}

%% TODO Titel: Ich-kann-Satz fehlt in der Bank (Platzhalter: Kettenname) – Nr. 13
%% TODO Anweisung: ein Satz über der ersten Teilaufgabe fehlt in der Bank – Nr. 13
\begin{aufgabe}{Ebene Figur: Flächeninhalt des Schattens eines geneigten Rechtecks bei senkrechtem Lichteinfall über eine Querschnittsskizze begründen}
\begin{teile}
\teil Ein rechteckiges Vordach ist gegen die Waagerechte geneigt, eine seiner Seiten verläuft waagerecht. Sonnenlicht fällt als parallele Strahlen senkrecht auf die Dachfläche und wirft einen rechteckigen Schatten auf den waagerechten Boden. Begründe mithilfe einer beschrifteten Querschnittsskizze, dass der Schatten einen größeren Flächeninhalt hat als das Vordach. \hfill (Abitur 2017 GK)
\end{teile}
\end{aufgabe}

%% TODO Titel: Ich-kann-Satz fehlt in der Bank (Platzhalter: Kettenname) – Nr. 14
%% TODO Anweisung: ein Satz über der ersten Teilaufgabe fehlt in der Bank – Nr. 14
\begin{aufgabe}{Vierecksflächen – Fehler finden, Begründen, Darstellungswechsel, Anwendung}
\begin{teile}
\teil Das Trapez $PQRS$ mit $P(1 | 0 | 2)$, $Q(11 | 0 | 2)$, $R(8 | 4 | 2)$, $S(4 | 4 | 2)$ hat die parallelen Seiten $PQ$ und $SR$. Mia rechnet: \rechnung{|QR| &= 5 \\ A &= \tfrac{1}{2} \cdot (10 + 4) \cdot 5 = 35} Finde den Fehler und rechne richtig.
\teil Begründe: Die Höhe eines Trapezes wird senkrecht zwischen den parallelen Seiten gemessen, nicht am Schenkel.
\teil Stelle für das Viereck $PQRS$ im Koordinatensystem einen Flächenterm „Rechteck plus rechtwinkliges Dreieck“ auf und berechne den Flächeninhalt.

\begin{ksys}[xmin=-1,xmax=7,ymin=-1,ymax=6] \punkt{0}{0}{P} \punkt{6}{0}{Q} \punkt{6}{5}{R} \punkt{0}{2}{S} \end{ksys}
\teil Eine schräge Zuschauertribüne ist ein symmetrisches Trapez mit den Ecken $P(0 | 0 | 0)$, $Q(20 | 0 | 0)$, $R(16 | 12 | 9)$, $S(4 | 12 | 9)$ (1 LE = 1 m). Pro Zuschauer rechnet man 0{,}6\,m². Wie viele Zuschauer passen auf die Tribüne? \leerfeld Zuschauer
\end{teile}
\end{aufgabe}

\clearpage
% Einheit 3 von 4 – Einheit 3 (bank/flaecheninhalt-und-volumen-im-raum/e3.jsonl, zusammenbau v0.1)
%% TODO Prüfungswort Einheit 3: nicht in der Marken-Zeile
%% TODO Zeitmarke Einheit 3: Marken-Zeile nicht lesbar
%% TODO Zweigzeile Teil 1: was hier gelernt wird (Satz in Schülersprache aus dem Einheitstitel) – Einheit 3
%% TODO Einheitentitel 3 nicht in der Mappe lesbar
\einheitenkopf[e3]{Einheit 3 von 4 · Einheit 3}
\setcounter{aufgabe}{14}

%% TODO Titel: Ich-kann-Satz fehlt in der Bank (Platzhalter: Kettenname) – Nr. 15
%% TODO Anweisung: ein Satz über der ersten Teilaufgabe fehlt in der Bank – Nr. 15
\begin{aufgabe}{Volumen und Oberfläche}
%% TODO \rechenplatz{n}: Schrittzahl des Grundfalls fehlt in der Bank (nur bei fester Schreibform, 2.3 b)
\begin{teile}
\teil Ein Zelt über einem quadratischen Boden, alle Dachkanten laufen in einer Spitze zusammen. Volumenformel? \\ \kreuz{mit Faktor ein Drittel} \\ \kreuz{ohne Faktor ein Drittel}
\teil Pyramide mit der Grundfläche $P(1 | 1 | 0)$, $Q(5 | 1 | 0)$, $R(5 | 5 | 0)$, $T(1 | 5 | 0)$ und der Spitze $S(3 | 3 | 6)$. Volumen? V = \leerfeld VE
\teil Das Dreieck $OPQ$ mit $P(0 | 3 | 4)$ und $Q(6 | 0 | 0)$ ist rechtwinklig in $O$. Die Pyramide $OPQS$ hat das Volumen 40, und es gilt $\overrightarrow{OS} \cdot \overrightarrow{OP} = 0$ und $\overrightarrow{OS} \cdot \overrightarrow{OQ} = 0$. Länge $|OS|$? \hfill (Abitur 2026 GK) \\ |OS| = \leerfeld
\teil Die Pyramide $PQRTS$ hat das Drachenviereck mit $P(0 | 0 | 0)$, $Q(3 | 2 | 0)$, $R(0 | 7 | 0)$, $T(-3 | 2 | 0)$ als Grundfläche und die Spitze $S(0 | 0 | 5)$. Volumen? \hfill (Abitur 2025 GK) \\ V = \leerfeld VE
\teil $P(1 | 1 | 0)$, $Q(5 | 1 | 0)$, $R(1 | 4 | 0)$, $P'(1 | 1 | 6)$, $Q'(5 | 1 | 6)$, $R'(1 | 4 | 6)$. Weise nach, dass $PQRP'Q'R'$ ein gerades Prisma ist, und berechne sein Volumen. V = \leerfeld VE
\teil Pyramide mit der quadratischen Grundfläche $P(-3 | -3 | 1)$, $Q(5 | -3 | 1)$, $R(5 | 5 | 1)$, $T(-3 | 5 | 1)$ und der Spitze $S(1 | 1 | 4)$. Oberflächeninhalt? \hfill (Abitur 2024 LK) \\ O = \leerfeld FE
\teil $P(1 | 0 | 2)$, $Q(4 | 4 | 2)$, $R(-4 | 10 | 2)$, $T(-7 | 6 | 2)$, $S(-1 | 5 | 8)$. Zeige, dass $PQRT$ ein Rechteck, aber kein Quadrat ist, und berechne das Volumen der Pyramide $PQRTS$.
\end{teile}
\end{aufgabe}

%% TODO Titel: Ich-kann-Satz fehlt in der Bank (Platzhalter: Kettenname) – Nr. 16
%% TODO Anweisung: ein Satz über der ersten Teilaufgabe fehlt in der Bank – Nr. 16
\begin{aufgabe}{Volumen und Oberfläche – Fehler finden, Begründen, Anwendung}
\begin{teile}
\teil Eine Pyramide hat die Grundfläche 24 und die Höhe 5. Lena rechnet: \rechnung{V &= 24 \cdot 5 \\ &= 120} Finde den Fehler und rechne richtig.
\teil Begründe: Eine Pyramide hat ein Drittel des Volumens eines Prismas mit gleicher Grundfläche und gleicher Höhe.
\teil Das Obergeschoss eines Veranstaltungssaals ist die Pyramide mit der waagerechten dreieckigen Grundfläche $P(0 | 0 | 12)$, $Q(0 | 24 | 12)$, $R(-20 | 6 | 12)$ und der Spitze $S(-8 | 8 | 30)$ (1 LE = 1 m). Die Lüftung braucht 0{,}5\,kW je 100\,m³. Reicht eine Anlage mit 8\,kW? \hfill (Abitur 2018 LK)
\end{teile}
\end{aufgabe}

\clearpage
% Einheit 4 von 4 – Einheit 4 (bank/flaecheninhalt-und-volumen-im-raum/e4.jsonl, zusammenbau v0.1)
%% TODO Prüfungswort Einheit 4: nicht in der Marken-Zeile
%% TODO Zeitmarke Einheit 4: Marken-Zeile nicht lesbar
%% TODO Zweigzeile Teil 1: was hier gelernt wird (Satz in Schülersprache aus dem Einheitstitel) – Einheit 4
%% TODO Einheitentitel 4 nicht in der Mappe lesbar
\einheitenkopf[e4]{Einheit 4 von 4 · Einheit 4}
\setcounter{aufgabe}{16}

%% TODO Titel: Ich-kann-Satz fehlt in der Bank (Platzhalter: Kettenname) – Nr. 17
%% TODO Anweisung: ein Satz über der ersten Teilaufgabe fehlt in der Bank – Nr. 17
\begin{aufgabe}{Zusammengesetzt und Verhältnisse}
%% TODO \rechenplatz{n}: Schrittzahl des Grundfalls fehlt in der Bank (nur bei fester Schreibform, 2.3 b)
\begin{teile}
\teil Ein Haus: Quader mit aufgesetztem Satteldach. Wie fasst du sein Volumen? \\ \kreuz{direkt} \\ \kreuz{als Summe} \\ \kreuz{als Differenz}
\teil Ein Quader hat die Grundfläche $PQRU$ mit $P(1 | 1 | 0)$, $Q(7 | 1 | 0)$, $R(7 | 7 | 0)$, $U(1 | 7 | 0)$ und die Höhe 4. Die Ebene durch $Q$, $U$ und $T(1 | 1 | 8)$ zerlegt den Quader. Berechne das Volumen des Teilkörpers, der $P$ enthält, und erläutere den Ansatz. V = \leerfeld VE
\teil Ein Würfel mit der Kantenlänge 6 hat eine Ecke im Ursprung, seine Kanten liegen auf den positiven Achsen. Die Ebene durch die Würfelkante auf der $x$-Achse und den Mittelpunkt $M(6 | 3 | 6)$ einer oberen Kante teilt den Würfel. Volumen des kleineren Teilkörpers? \hfill (Abitur 2025 LK) \\ V = \leerfeld VE
\teil Ein Sockel hat die Form eines Stumpfs einer geraden quadratischen Pyramide: Grundfläche mit der Seite 6\,m, Deckfläche mit der Seite 3\,m, Höhe 6\,m. Sein Volumen wird mit dem Term $\tfrac{1}{3} \cdot 36 \cdot 12 - \tfrac{1}{3} \cdot 9 \cdot 6$ berechnet. Begründe, dass der Term das Volumen liefert. \hfill (Abitur 2022 GK)
\teil Ein Quader und eine Pyramide haben dieselbe Grundfläche, die Spitze der Pyramide liegt in der Deckfläche des Quaders. Um wie viel Prozent ist das Volumen des Quaders größer als das der Pyramide? Rechne ohne konkrete Volumina. \hfill (Abitur 2021 GK) \\ \leerfeld \%
\teil Das Dreieck mit den Ecken $(0 | 0)$, $(8 | 0)$ und $(8 | 6)$ rotiert um die $x$-Achse. Oberflächeninhalt des entstehenden Körpers? \hfill (Abitur 2018 GK) \\ O = \leerfeld FE
\teil Eine zylindrische Vorratsdose mit dem Durchmesser 20\,cm fasst 10 Liter. Passt sie in ein Regalfach mit 30\,cm lichter Höhe? \hfill (Abitur 2022 GK)
\teil Die Pyramide mit den Ecken $O$, $P(6 | 0 | 0)$, $Q(0 | 7 | 0)$ und $R(0 | 0 | 5)$ wird von einer Ebene parallel zur $y$-$z$-Ebene in zwei volumengleiche Teile zerlegt. An welcher Stelle $t$ schneidet die Ebene die $x$-Achse? \hfill (Abitur 2022 LK) \\ t = \leerfeld
\end{teile}
\end{aufgabe}

%% TODO Titel: Ich-kann-Satz fehlt in der Bank (Platzhalter: Kettenname) – Nr. 18
%% TODO Anweisung: ein Satz über der ersten Teilaufgabe fehlt in der Bank – Nr. 18
\begin{aufgabe}{Zusammengesetzt und Verhältnisse – Fehler finden, Begründen, Darstellungswechsel, Anwendung}
\begin{teile}
\teil Ein Quader hat eine quadratische Grundfläche mit der Seite 6 und die Höhe 3. Die Ebene durch zwei Nachbarecken $Q$ und $U$ der Grundfläche und den Punkt $T$ senkrecht über ihrer gemeinsamen Nachbarecke $P$ in der Höhe 9 zerlegt den Quader. Paul rechnet für den Teilkörper mit $P$: \rechnung{V &= \tfrac{1}{3} \cdot \tfrac{1}{2} \cdot 6 \cdot 6 \cdot 9 = 54} Finde den Fehler und rechne richtig.
\teil Begründe: Schneidet eine Ebene parallel zur Grundfläche eine Teilpyramide ab, so ist ihr Volumenanteil der Kubus des Streckfaktors.
\teil Das Volumen eines Restkörpers ist $V(t) = 40 - 5t$ für $0 \le t \le 6$. Zeichne den Graphen von $V$ in das Koordinatensystem.

\begin{ksys}[xmin=0,xmax=7,ymin=0,ymax=45,xstep=1,ystep=5,xlabel=t,ylabel=V] \end{ksys}
\teil Ein Betonsockel ist der Stumpf einer geraden quadratischen Pyramide: unten Seite 1{,}2\,m, oben Seite 0{,}8\,m, Höhe 6\,m; die verlängerten Seitenkanten treffen sich 18\,m über dem Boden. Beton wiegt 2{,}4\,t je m³. Masse des Sockels? \hfill (Abitur 2018 LK) \\ \leerfeld[t]
\end{teile}
\end{aufgabe}

% Abhakseite „Das kann ich“ – Zone nur im Gesamt (\mitzone)
%% TODO Ich-kann-Sätze: Platzhalter wie in den Aufgabendateien
\begin{abhakseite}
\ifdefined\mitzone
\abhakgruppe{Kennst du schon}
\abhak{1}{Flächenformeln der ebenen Figuren (Dreieck, Trapez, Parallelogramm, Raute/Drachen, Kreis) kennen und anwenden}
\abhak{2}{Volumen- und Oberflächenformeln der Elementarkörper (Prisma, Pyramide, Zylinder, Kegel) kennen}
\abhak{3}{Verbindungsvektoren bilden, Beträge berechnen, Mittelpunkte bestimmen}
\abhak{4}{Skalarprodukt und „null heißt senkrecht“}
\abhak{5}{Satz des Pythagoras für Seitenhöhen und fehlende Katheten, Strahlensätze und Ähnlichkeit für Streckfaktoren}
\abhak{6}{Gleichungen aus Formeln lösen (nach einer Größe umstellen, Wurzelgleichungen, quadratische Terme, Prozentrechnung für Anteile)}
\abhak{7}{Flächenformeln der ebenen Figuren (Dreieck, Trapez, Parallelogramm, Raute/Drachen, Kreis) kennen und anwenden – Fehler finden}
\abhak{8}{Flächenformeln der ebenen Figuren (Dreieck, Trapez, Parallelogramm, Raute/Drachen, Kreis) kennen und anwenden – selbst rechnen}
\fi
\abhakgruppe{Einheit 1 · Dreiecksflächen}
\abhak{9}{Welche Figur, welche Formel?}
\abhak{10}{Dreiecksflächen}
\abhak{11}{Dreiecksflächen – Fehler finden, Begründen, Darstellungswechsel, Anwendung}
\abhakgruppe{Einheit 2 · Vierecksflächen}
\abhak{12}{Vierecksflächen}
\abhak{13}{Ebene Figur: Flächeninhalt des Schattens eines geneigten Rechtecks bei senkrechtem Lichteinfall über eine Querschnittsskizze begründen}
\abhak{14}{Vierecksflächen – Fehler finden, Begründen, Darstellungswechsel, Anwendung}
\abhakgruppe{Einheit 3 · Einheit 3}
\abhak{15}{Volumen und Oberfläche}
\abhak{16}{Volumen und Oberfläche – Fehler finden, Begründen, Anwendung}
\abhakgruppe{Einheit 4 · Einheit 4}
\abhak{17}{Zusammengesetzt und Verhältnisse}
\abhak{18}{Zusammengesetzt und Verhältnisse – Fehler finden, Begründen, Darstellungswechsel, Anwendung}
\end{abhakseite}

\end{document}
